% Compiled with LuaLaTeX. Main font is Roboto, loaded from the system TTFs via
% fontspec/fontconfig (the roboto LaTeX package is not installed here).
% Layout based on the Rover Resume base template
% (https://github.com/subidit/rover-resume): ruled section headers and a
% strict single-level bullet hierarchy, with run-in project titles that each
% collect their supporting artifacts as inline links.
\documentclass[11pt,a4paper]{article}

\usepackage[a4paper,margin=0.72in]{geometry}
\usepackage{fontspec}
\setmainfont{Roboto}[
  Ligatures=TeX,
  BoldFont={Roboto Medium},
  ItalicFont={Roboto Italic},
  BoldItalicFont={Roboto Medium Italic},
]
\usepackage[dvipsnames]{xcolor}
\definecolor{linkblue}{HTML}{1155CC}
\usepackage[colorlinks=true,urlcolor=linkblue,linkcolor=linkblue]{hyperref}
\usepackage{titlesec}
\usepackage{enumitem}

\setcounter{secnumdepth}{0}            % no section numbering
\titleformat{\section}{\large\bfseries\scshape}{}{0em}{}[\titlerule]
\titlespacing*{\section}{0pt}{8pt}{4pt}
\setlist[itemize]{topsep=1pt,itemsep=0pt,parsep=0pt,leftmargin=1.4em,label=\textbullet}
\setlength{\parindent}{0pt}
\pagestyle{empty}

% Run-in list of supporting artifacts after a project title.
\newcommand{\artifacts}[1]{\hspace{0.45em}{\footnotesize #1}}

% Projects are indented under their employer to establish visual hierarchy:
% the employer line is the flush-left anchor, projects nest one step in.
\newlength{\projindent}\setlength{\projindent}{1.3em}
\newlength{\projbulletindent}\setlength{\projbulletindent}{2.7em}

\begin{document}

% ===== Header =====
\begin{center}
  {\LARGE\bfseries Nikita Kazeev}\\[3pt]
  {\small
    \href{mailto:kazeevn@gmail.com}{kazeevn@gmail.com} \textperiodcentered{}
    \href{https://kazeevn.github.io}{Website} \textperiodcentered{}
    \href{https://github.com/kazeevn}{GitHub} \textperiodcentered{}
    \href{https://linkedin.com/in/nikita-kazeev}{LinkedIn} \textperiodcentered{}
    \href{https://scholar.google.com/citations?user=vamy2okAAAAJ}{Google Scholar} \textperiodcentered{}
    \href{https://orcid.org/0000-0002-5699-7634}{ORCID}%
  }\\[5pt]
  \textit{ Research Scientist at the intersection of AI and Physics }
\end{center}

% ===== PAGE 1: Work Experience & Projects =====
\section{Work Experience}
\textsc{\textbf{Constructor Group}}\quad{\small\color{black!55}\itshape Machine Learning Consultant} \hfill\textbf{2021--present}\par
\smallskip

\textsc{\textbf{National University of Singapore}}\quad{\small\color{black!55}\itshape Postdoc under \href{https://en.wikipedia.org/wiki/Konstantin\_Novoselov}{Kostya Novoselov}} \hfill\textbf{2022--present}\par
\vspace{2pt}
{\leftskip\projindent\textbf{Wyckoff Transformer for generating symmetric crystals}\artifacts{ \href{https://icml.cc/virtual/2025/poster/44595}{ICML 2025} \textperiodcentered{} \href{https://github.com/SymmetryAdvantage/WyckoffTransformer}{GitHub} }\par}
\begin{itemize}[leftmargin=\projbulletindent]
  \item Developed a coordinate-free permutation-invariant autoregressive Transformer encoder model for material generation and property prediction based on symmetry inductive bias
  \item Increased materials discovery yield (S.U.N.) by 1.24x
  \item Accelerated inference by 4 orders of magnitude, to 0.05 GPU ms per structure
  \item Implemented the model in PyTorch; production-ready package published on \href{https://pypi.org/project/wyckoff-transformer/}{PyPI}
  \item Orchestrated large-scale (10k structures) DFT computations with VASP at an HPC cluster
  \item Led a team of 6
\end{itemize}
\vspace{2pt}
{\leftskip\projindent\textbf{ML for defects in 2D crystals}\artifacts{ \href{https://www.nature.com/articles/s41699-023-00369-1}{npj 2D Mater. Appl.} \textperiodcentered{} \href{https://www.nature.com/articles/s41524-023-01062-z}{npj Comput. Mater.} \textperiodcentered{} \href{https://scholar.google.com/citations?view\_op=view\_citation&hl=en&user=vamy2okAAAAJ&sortby=pubdate&citation\_for\_view=vamy2okAAAAJ:uAPFzskPt0AC}{patent} \textperiodcentered{} \href{https://github.com/HSE-LAMBDA/ai4material\_design}{GitHub} }\par}
\begin{itemize}[leftmargin=\projbulletindent]
  \item Conceived a novel sparse representation of crystals with defects
  \item Achieved a 3.7x energy prediction error reduction
  \item Optimized training and inference: 4x less memory, 8x fewer GPU operations
  \item Developed code for reproducible parallel running of machine learning experiments
  \item Led a team of 3
\end{itemize}
\smallskip

\textsc{\textbf{CERN [Yandex $\rightarrow$ HSE University]}}\quad{\small\color{black!55}\itshape Researcher} \hfill\textbf{2014--2022}\par
{\small \href{https://breakthroughprize.org/Laureates/1/L3995}{2025 Breakthrough Prize in Fundamental Physics} as a member of LHCb. }\par
\vspace{2pt}
{\leftskip\projindent\textbf{Generative models for fast simulation}\artifacts{ \href{https://www.sciencedirect.com/science/article/pii/S0168900219300701}{paper 1} \textperiodcentered{} \href{https://arxiv.org/abs/1905.11825}{paper 2} \textperiodcentered{} \href{https://ml4physicalsciences.github.io/2019/files/NeurIPS\_ML4PS\_2019\_40.pdf}{paper 3} \textperiodcentered{} \href{https://github.com/yandexdataschool/QuantileTransformerTF}{code} }\par}
\begin{itemize}[leftmargin=\projbulletindent]
  \item Developed a GAN-based machine learning model for high-fidelity simulation of a Cherenkov detector trained on tabular data.
  \item Achieved approximately $10^5$ speed-up compared to the ab-initio simulator.
\end{itemize}
\vspace{2pt}
{\leftskip\projindent\textbf{Machine learning on noisy data}\artifacts{ \href{http://doi.org/10.1088/1748-0221/14/08/P08020}{paper 1} \textperiodcentered{} \href{https://ml4physicalsciences.github.io/2019/files/NeurIPS\_ML4PS\_2019\_122.pdf}{paper 2} \textperiodcentered{} \href{https://github.com/yandexdataschool/ML-sWeights-experiments/}{code} }\par}
\begin{itemize}[leftmargin=\projbulletindent]
  \item Developed a rigorous way to train ML algorithms on data with the label-noise model common in high-energy physics.
\end{itemize}
\vspace{2pt}
{\leftskip\projindent\textbf{Muon identification at the LHCb experiment at CERN}\artifacts{ \href{https://iopscience.iop.org/article/10.1088/1742-6596/1525/1/012100/pdf}{paper} \textperiodcentered{} \href{https://indico.cern.ch/event/708041/contributions/3272104/}{poster} \textperiodcentered{} \href{https://web.archive.org/web/20260209172339/https://idao.world/history/\#idao-2019}{IDAO-2019} }\par}
\begin{itemize}[leftmargin=\projbulletindent]
  \item Solved a classification problem over tabular data with noisy labels under timing constraints.
  \item Developed a gradient boosting model that reduced the false positive rate by 30\% in the critical low-track-momentum region.
  \item Integrated the model into LHCb production, mostly in C++.
  \item Packaged the work as a data science competition problem for IDAO-2019.
\end{itemize}
\smallskip

\newpage

% ===== PAGE 2: Education & Activities =====
\section{Education}
\textbf{HSE University \& Sapienza Università di Roma}\hfill\textbf{2016--2020}\par
\textit{PhD in Computer Science and Physics (Double Degree)}\par
{\leftskip\projindent Supervisors: \href{https://www.linkedin.com/in/andreyustyuzhanin/}{Andrey Ustyuzhanin} \& \href{https://it.linkedin.com/in/barbara-sciascia-234935a}{Barbara Sciascia}\par}
{\leftskip\projindent Thesis: \href{https://repository.cern/records/034yx-nt191}{Machine Learning for particle identification in the LHCb detector}\par}
\smallskip

\textbf{Product Management course at Yandex}\hfill\textbf{2018}\par
\textit{Focused coursework in product strategy and execution.}\par
\smallskip

\textbf{Yandex School of Data Analysis}\hfill\textbf{2013--2015}\par
\textit{Master's-level CS programme}\par
{\leftskip\projindent Algorithms, machine learning, deep learning, and distributed systems.\par}
\smallskip

\textbf{Moscow Institute of Physics and Technology}\par
\textit{MS in Physics}\hfill\textbf{2014--2016}\par
{\leftskip\projindent Optimisation of data processing of the LHCb experiment.\par}
\textit{BS in Physics}\hfill\textbf{2010--2014}\par
{\leftskip\projindent Wave-packet molecular dynamics of electrons in nonideal plasma.\par}
\smallskip

\textbf{International Junior Science Olympiad (IJSO)}\hfill\textbf{Korea, 2008}\par
{\leftskip\projindent Silver medal\par}
\smallskip


\section{Mentorship}
\begin{itemize}
  \item Mentored 9 students (\href{https://www.hse.ru/edu/vkr/219430466}{1}, \href{https://www.hse.ru/edu/vkr/296294066}{2}, \href{https://www.hse.ru/edu/vkr/219524130}{3}, \href{https://www.hse.ru/en/edu/vkr/470561417}{4}, \href{https://www.hse.ru/en/edu/vkr/470935224}{5}, \href{https://www.hse.ru/ma/datasci/students/diplomas/631565328}{6}, \href{https://www.hse.ru/ma/datasci/students/diplomas/631565422}{7}, \href{https://www.hse.ru/edu/vkr/624890096}{8}, \href{https://www.hse.ru/ba/ami/students/diplomas/834351454}{9}), 3 interns, and student workshop projects (\href{https://sochisirius.ru/news/3148}{1}, \href{https://web.archive.org/web/20220512040654/https://academy.yandex.ru/posts/vossozdat-arkhitekturu-bazy-dannykh-i-nauchit-golosovogo-pomoschnika-predskazyvat-namereniya-polzovatelya}{2}).
  \item Co-PI of a \$3.4m AI Singapore grant on multiscale machine learning.
\end{itemize}
\section{Teaching \& Outreach}
\begin{itemize}
  \item Pitch research to the general public through major newspaper coverage (\href{https://www.kommersant.ru/doc/7832307}{1}, \href{https://www.kommersant.ru/doc/6122817}{2}), blog posts (\href{https://ifim.nus.edu.sg/tailor-made-2d-materials-became-closer-with-the-new-machine-learning-algorithm-developed-at-nus/}{1}, \href{https://communities.springernature.com/posts/sparse-representation-for-machine-learning-the-properties-of-defects-in-2d-materials}{2}), a \href{https://elementy.ru/events/435506/IV\_Nauchnye\_boi\_NIU\_VShE\_na\_Dne\_Vyshki}{science slam}, interviews (\href{https://web.archive.org/web/20210518043835/https://academy.yandex.ru/posts/beresh-i-razbiraeshsya}{1}, \href{https://habr.com/en/company/yandex/blog/422761/}{2}), and a public \href{https://www.youtube.com/watch?v=zLFZ\_z-7-JI&feature=youtu.be&t=7779}{talk}.
  \item Taught at the Machine Learning in High Energy Physics Summer School in \href{https://indico.cern.ch/event/439520/}{2015}, \href{https://indico.cern.ch/event/497368/}{2016}, \href{https://indico.cern.ch/event/613571/}{2017}, \href{https://indico.cern.ch/event/687473/}{2018}, \href{https://indico.cern.ch/event/768915/}{2019}, \href{https://indico.cern.ch/event/838377/}{2020}, \href{https://indico.cern.ch/event/1025052/}{2021}, and \href{https://indico.cern.ch/event/1229514/}{2023}.
  \item Lecturer in the \href{https://www.hse.ru/en/news/edu/259378493.html}{award-winning} Coursera Advanced Machine Learning \href{https://web.archive.org/web/20220127183223/https://www.coursera.org/specializations/aml}{  specialisation}
  \item Delivered more than 20 \href{https://docs.google.com/document/d/1IOlBse4q0YgEoL\_RCJTvJfbFhFaD\_YMAIJJ3W77B\_X0/edit?usp=sharing}{conference talks}.
\end{itemize}
\section{Service}
\begin{itemize}
  \item Reviewer for NeurIPS, RSC Advances, Machine Learning: Science and Technology, Digital Discovery
  \item Peer Staff Supporter at NUS, serving as a first-line support for mental wellbeing.
\end{itemize}

\newpage

% ===== PAGE 3: Skills & Publications =====
\section{Skills}
\begin{itemize}
  \item \textbf{ML \& AI} --- Structured, non-text, domain-specific data
  \item \textbf{Python} --- Research coding
  \item \textbf{PyTorch} --- Deep learning
  \item \textbf{C++} --- Production
  \item \textbf{Linux} --- Administration
  \item \textbf{Research Writing \& Speaking}
  \item \textbf{Physics} --- Deep domain understanding
  \item \textbf{Leadership \& Mentorship}
\end{itemize}

\section{Selected Publications}
\begin{enumerate}[leftmargin=1.5em,topsep=1pt,itemsep=1pt,parsep=0pt]
  \item \textbf{Kazeev, Nikita}, and Ian Babich. \href{https://openreview.net/forum?id=GhUK6VGW67}{Foresight-Phys: A Benchmark for Forecasting the Results of Physical Experiments}. Forecasting as a New Frontier of Intelligence Workshop at ICML 2026
  \item Artem Dembitskiy, Shuya Yamazaki, Artem Maevskiy, \textbf{Nikita Kazeev}, Roman A. Eremin, Semen Budennyy, Kedar Hippalgaonkar, Antonio Helio Castro Neto, and Andrey E. Ustyuzhanin. \href{https://openreview.net/forum?id=xFI6K4ZQDp}{Generative Pipeline for Discovering Solid-State Battery Materials with Universal Atomistic Potentials}. AI for Science Workshop at ICML 2026
  \item Nian Liu, \textbf{Nikita Kazeev}, Stephen Gregory Dale, Artem Maevskiy, Yuwei Zeng, Ryoji Kubo, Pengru Huang, Thomas Laurent, Yann LeCun, Kostya S. Novoselov, and Xavier Bresson. \href{https://arxiv.org/abs/2605.14759}{Crys-JEPA: Accelerating Crystal Discovery via Embedding Screening and Generative Refinement}. arXiv preprint arXiv:2605.14759, May 2026.
  \item Alexandre Duval, Siddharth Betala, Samuel P. Gleason, Andy Xu, Georgia Channing, Daniel Levy, Ali Ramlaoui, Clémentine Fourrier, Chaitanya K. Joshi, \textbf{Nikita Kazeev}, Sékou-Oumar Kaba, Félix Therrien, Alex Hernández-García, Rocío Mercado, N M Anoop Krishnan. \href{https://arxiv.org/abs/2512.04562}{LeMat-GenBench: Bridging the gap between crystal generation and materials discovery}. AI for Accelerated Materials Design - NeurIPS 2025 / arXiv preprint arXiv:2512.04562, October 2025.
  \item \textbf{Kazeev, Nikita}, et al. \href{https://arxiv.org/abs/2503.02407}{WyckoffTransformer: Generation of Symmetric Crystals}. ICML 2025, May 2025.
  \item \textbf{Kazeev, N.}, Al-Maeeni, A.R., Romanov, I. et al. \href{https://doi.org/10.1038/s41524-023-01062-z}{Sparse representation for machine learning the properties of defects in 2D materials}. npj Comput Mater 9, 113 (2023).
  \item Derkach, D., \textbf{Kazeev, N.}, Ratnikov, F., Ustyuzhanin, A., \& Volokhova, A. (2019). \href{https://doi.org/10.1016/j.nima.2019.01.031}{Cherenkov detectors fast simulation using neural networks}. \textit{Nuclear Instruments and Methods in Physics Research Section A}. \textbf{Alphabetic order, I'm the corresponding author.}
\end{enumerate}

\end{document}